\documentclass[11pt,oneside]{article}
\pdfminorversion=7
\usepackage{QuizStyle}
\hypersetup{
 pdftitle={20261224 Final for MATH312 Analysis II},
 pdfauthor={Jungwoo Kim},
 pdfsubject={Rudin Chapters 9 and 11 exercises and solutions for MATH312 Analysis II},
 pdfcreator={LaTeX}
}
\begin{document}
\quiztitle{20261224 Final}

\begin{exercise}{9.1}
If $S$ is a nonempty subset of a vector space $X$, prove (as asserted in Section 9.1) that the span of $S$ is a vector space.
\end{exercise}
\begin{solution}
By definition, $\Span(S)$ consists of all finite linear combinations of elements of $S$. In particular, $\mathbf 0=0\mathbf s\in\Span(S)$ for any $\mathbf s\in S$. If
\[
 \mathbf x=\sum_{i=1}^k a_i\mathbf s_i,\qquad
 \mathbf y=\sum_{j=1}^{\ell}b_j\mathbf t_j
\]
are in $\Span(S)$, then $\mathbf x+\mathbf y$ is again a finite linear combination of elements of $S$. For every scalar $c$, so is $c\mathbf x=\sum_{i=1}^k(ca_i)\mathbf s_i$. Thus $\Span(S)$ contains $\mathbf 0$ and is closed under addition and scalar multiplication. It is therefore a vector subspace of $X$.
\end{solution}

\begin{exercise}{9.2}
Prove (as asserted in Section 9.6) that $\mathbf B\mathbf A$ is linear if $\mathbf A$ and $\mathbf B$ are linear transformations. Prove also that $\mathbf A^{-1}$ is linear and invertible whenever $\mathbf A$ is an invertible linear transformation.
\end{exercise}
\begin{solution}
For vectors $\mathbf x,\mathbf y$ and scalars $c,d$, linearity gives
\[
 (\mathbf B\mathbf A)(c\mathbf x+d\mathbf y)
 =\mathbf B(c\mathbf A\mathbf x+d\mathbf A\mathbf y)
 =c\mathbf B\mathbf A\mathbf x+d\mathbf B\mathbf A\mathbf y.
\]
If $\mathbf A$ is invertible, let $\mathbf u=\mathbf A\mathbf x$ and $\mathbf v=\mathbf A\mathbf y$. Then
\[
 \mathbf A^{-1}(c\mathbf u+d\mathbf v)
 =\mathbf A^{-1}\bigl(\mathbf A(c\mathbf x+d\mathbf y)\bigr)
 =c\mathbf A^{-1}\mathbf u+d\mathbf A^{-1}\mathbf v.
\]
Finally, the inverse of $\mathbf A^{-1}$ is $\mathbf A$.
\end{solution}

\begin{exercise}{9.3}
Assume $\mathbf A\colon X\to Y$ is linear and $\mathbf A\mathbf x=\mathbf 0$ only when $\mathbf x=\mathbf 0$. Prove that $\mathbf A$ is one-to-one.
\end{exercise}
\begin{solution}
If $\mathbf A\mathbf x=\mathbf A\mathbf y$, linearity yields $\mathbf A(\mathbf x-\mathbf y)=\mathbf 0$. By hypothesis, $\mathbf x-\mathbf y=\mathbf 0$, so $\mathbf x=\mathbf y$. Thus $\mathbf A$ is injective.
\end{solution}

\begin{exercise}{9.4}
Prove (as asserted in Section 9.30) that the null spaces and ranges of linear transformations are vector spaces.
\end{exercise}
\begin{solution}
Let $\mathbf A\colon X\to Y$ be linear. Its null space $\ker\mathbf A$ contains $\mathbf 0$ and, for $\mathbf x,\mathbf y\in\ker\mathbf A$,
\[
 \mathbf A(c\mathbf x+d\mathbf y)=c\mathbf A\mathbf x+d\mathbf A\mathbf y=\mathbf 0.
\]
Hence $\ker\mathbf A$ is a subspace of $X$. The range $\im\mathbf A$ contains $\mathbf 0$; if $\mathbf u=\mathbf A\mathbf x$ and $\mathbf v=\mathbf A\mathbf y$, then $c\mathbf u+d\mathbf v=\mathbf A(c\mathbf x+d\mathbf y)\in\im\mathbf A$. Thus $\im\mathbf A$ is a subspace of $Y$.
\end{solution}

\begin{exercise}{9.5}
Prove that to every $\mathbf A\in\mathcal L(\R^n,\R)$ corresponds a unique $\mathbf y\in\R^n$ such that $\mathbf A\mathbf x=\mathbf x\cdot\mathbf y$. Prove also that $\|\mathbf A\|=\|\mathbf y\|$.

\textit{Hint.} Under certain conditions, equality holds in the Schwarz inequality.
\end{exercise}
\begin{solution}
Let $\mathbf e_1,\ldots,\mathbf e_n$ be the standard basis and put $y_i=\mathbf A\mathbf e_i$. Since $\mathbf x=\sum_i x_i\mathbf e_i$,
\[
 \mathbf A\mathbf x=\sum_{i=1}^n x_i y_i=\mathbf x\cdot\mathbf y.
\]
Taking $\mathbf x=\mathbf e_i$ shows that each $y_i$ is uniquely determined. By the Cauchy--Schwarz inequality, $|\mathbf A\mathbf x|\le\|\mathbf y\|\,\|\mathbf x\|$, so $\|\mathbf A\|\le\|\mathbf y\|$. If $\mathbf y\ne\mathbf 0$, choose $\mathbf x=\mathbf y/\|\mathbf y\|$; then $|\mathbf A\mathbf x|=\|\mathbf y\|$. If $\mathbf y=\mathbf 0$, both norms vanish. Hence $\|\mathbf A\|=\|\mathbf y\|$.
\end{solution}

\begin{exercise}{9.6}
Define $f(0,0)=0$ and
\[
 f(x,y)=\frac{xy}{x^2+y^2}\qquad((x,y)\ne(0,0)).
\]
Prove that $\partial_1f$ and $\partial_2f$ exist at every point of $\R^2$, although $f$ is not continuous at $(0,0)$.
\end{exercise}
\begin{solution}
Away from the origin, $f$ is a quotient of differentiable functions with nonzero denominator. At the origin,
\[
 \partial_1f(0,0)=\lim_{h\to0}\frac{f(h,0)-f(0,0)}h=0,
 \quad
 \partial_2f(0,0)=\lim_{h\to0}\frac{f(0,h)-f(0,0)}h=0.
\]
However, $f(t,t)=1/2$ for every $t\ne0$, whereas $f(0,0)=0$. Thus $f$ is not continuous at the origin.
\end{solution}

\begin{exercise}{9.7}
Suppose that $f$ is a real-valued function defined in an open set $E\subseteq\R^n$, and that $\partial_1f,\ldots,\partial_nf$ are bounded in $E$. Prove that $f$ is continuous in $E$.

\textit{Hint.} Proceed as in the proof of Theorem 9.21.
\end{exercise}
\begin{solution}
Fix $\mathbf a\in E$ and choose a ball $B(\mathbf a,r)\subseteq E$. Let $M$ bound the absolute values of all partial derivatives on $E$. For $\mathbf h$ sufficiently small, successively change the coordinates of $\mathbf a$ to those of $\mathbf a+\mathbf h$; every resulting line segment stays in the ball. Applying the one-variable mean value theorem on each segment and adding gives
\[
 |f(\mathbf a+\mathbf h)-f(\mathbf a)|
 \le M\sum_{i=1}^n|h_i|
 \le M\sqrt n\,\|\mathbf h\|\longrightarrow0.
\]
Hence $f$ is continuous at $\mathbf a$, and therefore on $E$.
\end{solution}

\begin{exercise}{9.8}
Suppose that $f$ is a differentiable real function in an open set $E\subseteq\R^n$, and that $f$ has a local maximum at $\mathbf a\in E$. Prove that $f'(\mathbf a)=0$.
\end{exercise}
\begin{solution}
For each standard basis vector $\mathbf e_i$, the function $t\mapsto f(\mathbf a+t\mathbf e_i)$ has a local maximum at $t=0$. Its derivative there is therefore zero: $\partial_i f(\mathbf a)=0$. Since the derivative $f'(\mathbf a)$ is represented by these partial derivatives (Theorem 9.17), $f'(\mathbf a)=0$.
\end{solution}

\begin{exercise}{9.9}
If $\mathbf f$ is a differentiable mapping of a connected open set $E\subseteq\R^n$ into $\R^m$, and $\mathbf f'(\mathbf x)=\mathbf 0$ for every $\mathbf x\in E$, prove that $\mathbf f$ is constant in $E$.
\end{exercise}
\begin{solution}
Fix $\mathbf a\in E$. On any ball contained in $E$, the restriction of $\mathbf f$ to a line segment has derivative zero by the chain rule; the one-variable mean value theorem applied to each component shows that $\mathbf f$ is constant on the ball. Consequently, the set
\[
 A=\{\mathbf x\in E:\mathbf f(\mathbf x)=\mathbf f(\mathbf a)\}
\]
is open relative to $E$. Its complement is also open, by the same local constancy. Since $E$ is connected and $A$ is nonempty, $A=E$.
\end{solution}

\begin{exercise}{9.10}
If $f$ is a real function on a convex open set $E\subseteq\R^n$ such that $\partial_1f(\mathbf x)=0$ for every $\mathbf x\in E$, prove that $f(\mathbf x)$ depends only on $x_2,\ldots,x_n$.

Show that convexity can be replaced by a weaker condition, but that some condition is required. For example, for $n=2$ and a horseshoe-shaped $E$, the statement may fail.
\end{exercise}
\begin{solution}
Fix $x_2,\ldots,x_n$. The fiber
\[
 I=\{t\in\R:(t,x_2,\ldots,x_n)\in E\}
\]
is an interval when $E$ is convex. On $I$, the function $h(t)=f(t,x_2,\ldots,x_n)$ satisfies $h'(t)=0$, hence is constant by the one-variable mean value theorem. Thus two points of $E$ with the same last $n-1$ coordinates have the same $f$-value.

The argument only requires that every nonempty fiber $I$ be an interval; convexity is sufficient but not necessary. To see why a restriction is needed, in $\R^2$ take the connected open U-shaped set
\[
 E=((-2,-1)\times(-1,3))\cup((1,2)\times(-1,3))
   \cup((-2,2)\times(-1,0)).
\]
Set $f=0$ on the left arm and the bottom strip, and set $f(x,y)=\psi(y)$ on the right arm for $y\ge0$, where $\psi(y)=e^{-1/y^2}$ for $y>0$ and $\psi(y)=0$ for $y\le0$. The definitions agree where the sets overlap; $\partial_1f=0$ everywhere. But at any fixed $y>0$, $f$ is zero on the left arm and positive on the right arm, so it does not depend only on $y$.
\end{solution}

\begin{exercise}{9.11}
If $f$ and $g$ are differentiable real functions in $\R^n$, prove that
\[
 \nabla(fg)=f\nabla g+g\nabla f
\]
and that $\nabla(1/f)=-f^{-2}\nabla f$ wherever $f\ne0$.
\end{exercise}
\begin{solution}
Apply the one-variable product rule to each partial derivative:
$\partial_i(fg)=f\partial_i g+g\partial_i f$ for $1\le i\le n$. Collecting the coordinates proves the first identity. At points where $f\ne0$, the one-variable chain rule gives $\partial_i(1/f)=-f^{-2}\partial_i f$, proving the second identity.
\end{solution}

\begin{exercise}{9.12}
Fix $0<a<b$, and define $\mathbf f=(f_1,f_2,f_3)\colon\R^2\to\R^3$ by
\[
 \begin{aligned}
 f_1(s,t)&=(b+a\cos s)\cos t,\\
 f_2(s,t)&=(b+a\cos s)\sin t,\\
 f_3(s,t)&=a\sin s.
 \end{aligned}
\]
Describe the range $K$ of $\mathbf f$ (a compact subset of $\R^3$).
\begin{exercisesubparts}
\item Show that precisely four $\mathbf p\in K$ have the property that $\nabla f_1(s,t)=\mathbf0$ whenever $\mathbf f(s,t)=\mathbf p$, and find them.
\item Determine all $\mathbf q\in K$ such that $\nabla f_3(s,t)=\mathbf0$ whenever $\mathbf f(s,t)=\mathbf q$.
\item Classify the four points in (a) as a local maximum, a local minimum, or saddle points of $f_1$. Which points in (b) correspond to maxima and minima of $f_3$?
\item For irrational $\lambda$, let $\mathbf g(t)=\mathbf f(t,\lambda t)$. Prove that $\mathbf g$ is a one-to-one map of $\R$ onto a dense subset of $K$, and that
\[
 \|\mathbf g'(t)\|^2=a^2+\lambda^2(b+a\cos t)^2.
\]
\end{exercisesubparts}
\end{exercise}
\begin{solution}
Put $\rho=\sqrt{X^2+Y^2}$ for a point $(X,Y,Z)\in\R^3$. Since $b+a\cos s>0$, the parametrization gives
\[
 K=\{(X,Y,Z)\in\R^3:(\rho-b)^2+Z^2=a^2\}.
\]
This is the torus obtained by rotating the circle of radius $a$ centered at $(b,0)$ about the $Z$-axis; it is compact because it is the image of $[0,2\pi]^2$.

\begin{exercisesubparts}
\item The partial derivatives are
\[
 \partial_s f_1=-a\sin s\cos t,
 \qquad \partial_t f_1=-(b+a\cos s)\sin t.
\]
Thus $\sin t=0$ and then $\sin s=0$, so $s,t\in\pi\Z$. Their images are exactly
\[
 (b+a,0,0),\quad (-(b+a),0,0),\quad
 (b-a,0,0),\quad (-(b-a),0,0).
\]
Each point of $K$ has a unique pair of angles modulo $2\pi$, so there are no further points.
\item Since $\partial_s f_3=a\cos s$ and $\partial_t f_3=0$, the required points are the two circles
\[
 \{(b\cos t,b\sin t,a):t\in\R\},\qquad
 \{(b\cos t,b\sin t,-a):t\in\R\}.
\]
\item At the four critical pairs, the mixed second derivative of $f_1$ vanishes and
\[
 \partial_{ss}f_1=-a\cos s\cos t,
 \qquad \partial_{tt}f_1=-(b+a\cos s)\cos t.
\]
The point $(b+a,0,0)$ is the maximum and $(-(b+a),0,0)$ is the minimum; the other two have second derivatives of opposite signs and are saddles. The upper circle $Z=a$ consists of maxima of $f_3$, and the lower circle $Z=-a$ consists of minima.
\item If $\mathbf g(t_1)=\mathbf g(t_2)$, the uniqueness of torus angles modulo $2\pi$ implies
\[
 t_1-t_2=2\pi k,\qquad \lambda(t_1-t_2)=2\pi\ell
 \quad(k,\ell\in\Z).
\]
Since $\lambda$ is irrational, $k=\ell=0$, so $t_1=t_2$.
For a fixed first angle $s$, the values $t=s+2\pi k$ have second angles $\lambda s+2\pi k\lambda$ modulo $2\pi$. To see that the rotation orbit is dense, apply the pigeonhole principle to the fractional parts of $0,\lambda,\ldots,N\lambda$: two lie within $1/N$ of each other, so some nonzero integer multiple of $\lambda$ differs from an integer by a number $\delta$ with $0<|\delta|<1/N$. Changing its sign if necessary, successive integer multiples of the small positive step $\delta$ approximate every point modulo $1$ within $\delta$. Since $N$ is arbitrary, $\{k\lambda\bmod1:k\in\Z\}$ is dense. Thus every pair of angles, and hence every point of $K$, is approximated by points of $\mathbf g(\R)$.
Finally,
\[
 \|\partial_s\mathbf f\|^2=a^2,\quad
 \|\partial_t\mathbf f\|^2=(b+a\cos s)^2,\quad
 \partial_s\mathbf f\cdot\partial_t\mathbf f=0.
\]
The chain rule gives $\mathbf g'=\partial_s\mathbf f+\lambda\partial_t\mathbf f$ at $(s,t)=(t,\lambda t)$; the claimed identity follows.
\end{exercisesubparts}
\end{solution}

\begin{exercise}{9.13}
Suppose $\mathbf f\colon\R\to\R^3$ is differentiable and $\|\mathbf f(t)\|=1$ for every $t$. Prove that $\mathbf f'(t)\cdot\mathbf f(t)=0$, and interpret this geometrically.
\end{exercise}
\begin{solution}
Differentiate $\mathbf f(t)\cdot\mathbf f(t)=1$ using the product rule:
\[
 2\mathbf f'(t)\cdot\mathbf f(t)=0.
\]
Thus the velocity vector of a differentiable curve on the unit sphere is perpendicular to its radius vector; equivalently, it is tangent to the sphere.
\end{solution}

\begin{exercise}{9.14}
Define $f(0,0)=0$ and
\[
 f(x,y)=\frac{x^3}{x^2+y^2}\qquad((x,y)\ne(0,0)).
\]
\begin{exercisesubparts}
\item Prove that $\partial_1 f$ and $\partial_2 f$ are bounded functions on $\R^2$ (so $f$ is continuous).
\item For every unit vector $\mathbf u\in\R^2$, show that $D_{\mathbf u}f(0,0)$ exists and has absolute value at most $1$.
\item Let $\boldsymbol\gamma\colon\R\to\R^2$ be differentiable, with $\boldsymbol\gamma(0)=\mathbf0$ and $\|\boldsymbol\gamma'(0)\|>0$. Put $g(t)=f(\boldsymbol\gamma(t))$. Prove that $g$ is differentiable for every $t\in\R$, and that $g\in C^1$ whenever $\boldsymbol\gamma\in C^1$.
\item Prove nonetheless that $f$ is not differentiable at $(0,0)$. 
\textit{Hint.} Formula (40) fails.
\end{exercisesubparts}
\end{exercise}
\begin{solution}
\begin{exercisesubparts}
\item For $(x,y)\ne(0,0)$, direct differentiation gives
\[
 \partial_1 f=\frac{x^2(x^2+3y^2)}{(x^2+y^2)^2},\qquad
 \partial_2 f=-\frac{2x^3y}{(x^2+y^2)^2}.
\]
Both have absolute value at most $3$ and $2$, respectively. At the origin, the coordinate difference quotients give $\partial_1f(0,0)=1$, $\partial_2f(0,0)=0$; thus both partial derivatives are bounded everywhere. By Exercise 9.7, $f$ is continuous.
\item For $\mathbf u=(u_1,u_2)$ with $\|\mathbf u\|=1$, homogeneity gives
\[
 D_{\mathbf u}f(0,0)=\lim_{t\to0}\frac{f(t\mathbf u)}t=u_1^3,
\]
whose absolute value is at most $1$.
\item The function $f$ is homogeneous of degree one, satisfies $|f(\mathbf z)|\le\|\mathbf z\|$, and has bounded first partial derivatives off the origin. At any $t_0$ with $\boldsymbol\gamma(t_0)\ne\mathbf0$, the usual chain rule applies. If $\boldsymbol\gamma(t_0)=\mathbf0$, write
\[
 \boldsymbol\gamma(t_0+h)=h\mathbf v+o(|h|),
 \qquad\mathbf v=\boldsymbol\gamma'(t_0).
\]
The bounded partial derivatives imply that $f$ is locally Lipschitz (by the line-segment argument of Exercise 9.7). Consequently,
\[
 g(t_0+h)=f(h\mathbf v)+o(|h|)=h f(\mathbf v)+o(|h|),
\]
so $g'(t_0)=f(\mathbf v)$. This also covers $\mathbf v=\mathbf0$.
If $\boldsymbol\gamma\in C^1$, continuity of $g'$ is automatic where $\boldsymbol\gamma\ne\mathbf0$. At a zero $t_0$ with $\boldsymbol\gamma'(t_0)=\mathbf0$, boundedness of $\nabla f$ away from the origin gives $|g'(t)|\le C\|\boldsymbol\gamma'(t)\|\to0$. At a zero with $\boldsymbol\gamma'(t_0)=\mathbf v\ne\mathbf0$, the directions $\boldsymbol\gamma(t)/\|\boldsymbol\gamma(t)\|$ approach $\pm\mathbf v/\|\mathbf v\|$. As $\nabla f$ is homogeneous of degree zero and continuous off the origin,
\[
 g'(t)\to\nabla f(\mathbf v)\cdot\mathbf v=f(\mathbf v)=g'(t_0),
\]
where the penultimate equality follows by differentiating $f(t\mathbf v)=tf(\mathbf v)$. Hence $g\in C^1$.
\item Differentiability of $f$ at the origin would force
$D_{\mathbf u}f(0,0)=\partial_1f(0,0)u_1+\partial_2f(0,0)u_2=u_1$.
But (b) gives $D_{\mathbf u}f(0,0)=u_1^3$. For $\mathbf u=(1/\sqrt2,1/\sqrt2)$ these values differ. Thus $f$ is not differentiable there.
\end{exercisesubparts}
\end{solution}

\begin{exercise}{9.15}
Define $f(0,0)=0$ and, for $(x,y)\ne(0,0)$, put
\[
 f(x,y)=x^2+y^2-2x^2y-\frac{4x^6y^2}{(x^4+y^2)^2}.
\]
\begin{exercisesubparts}
\item Prove that $4x^4y^2\le(x^4+y^2)^2$ for all $(x,y)\in\R^2$, and conclude that $f$ is continuous.
\item For $0\le\theta\le2\pi$, define $g_\theta(t)=f(t\cos\theta,t\sin\theta)$. Show that $g_\theta(0)=g_\theta'(0)=0$ and $g_\theta''(0)=2$, so that each restriction to a line through the origin has a strict local minimum there.
\item Show that the origin is nevertheless not a local minimum, since $f(x,x^2)=-x^4$.
\end{exercisesubparts}
\end{exercise}
\begin{solution}
\begin{exercisesubparts}
\item The inequality follows from $(x^4-y^2)^2\ge0$. Thus
\[
 0\le\frac{4x^6y^2}{(x^4+y^2)^2}\le x^2
 \qquad((x,y)\ne(0,0)).
\]
All the remaining terms tend to zero at the origin, proving continuity there; continuity elsewhere is immediate.
\item Write $c=\cos\theta$, $s=\sin\theta$. Then
\[
 g_\theta(t)=t^2-2t^3c^2s-
 \frac{4t^8c^6s^2}{(t^4c^4+t^2s^2)^2}
 \quad(t\ne0).
\]
If $s\ne0$, the last term is $O(t^4)$ as $t\to0$; if $s=0$, it vanishes identically. In either case, $g_\theta$ is twice differentiable near $0$ (for $s\ne0$ the displayed expression extends smoothly across $t=0$), and
$g_\theta(t)=t^2+o(t^2)$. Therefore $g_\theta(0)=g_\theta'(0)=0$ and $g_\theta''(0)=2$. In particular, $g_\theta(t)>0$ for all sufficiently small nonzero $t$, with the neighborhood allowed to depend on $\theta$.
\item For $x\ne0$, substitute $y=x^2$:
\[
 f(x,x^2)=x^2+x^4-2x^4-
 \frac{4x^{10}}{(2x^4)^2}=-x^4<0=f(0,0).
\]
Therefore the origin is not a local minimum of $f$.
\end{exercisesubparts}
\end{solution}

\begin{exercise}{9.16}
Show that continuity of $f'$ at $a$ is needed in the inverse function theorem, even for $n=1$: let
\[
 f(t)=t+2t^2\sin(1/t)\quad(t\ne0),\qquad f(0)=0.
\]
Show that $f'(0)=1$, that $f'$ is bounded on $(-1,1)$, and that $f$ is not one-to-one in any neighborhood of $0$.
\end{exercise}
\begin{solution}
Since $f(t)/t=1+2t\sin(1/t)\to1$, we have $f'(0)=1$. For $t\ne0$,
\[
 f'(t)=1+4t\sin(1/t)-2\cos(1/t),
\]
which is bounded on $(-1,1)$. At $t_k=1/(2\pi k)$, $f'(t_k)=-1$; at $s_k=1/((2k+1)\pi)$, $f'(s_k)=3$. Both sequences tend to $0$. A continuous injective function on an interval must be strictly monotone, but a differentiable monotone function cannot have derivatives of opposite signs. Hence $f$ is not injective on any interval about $0$, despite $f'(0)\ne0$.
\end{solution}

\begin{exercise}{9.17}
Let $\mathbf f=(f_1,f_2)\colon\R^2\to\R^2$ be given by
\[
 f_1(x,y)=e^x\cos y,\qquad f_2(x,y)=e^x\sin y.
\]
\begin{exercisesubparts}
\item What is the range of $\mathbf f$?
\item Show that the Jacobian determinant never vanishes, so $\mathbf f$ is locally one-to-one everywhere, but is not one-to-one on $\R^2$.
\item Put $\mathbf a=(0,\pi/3)$ and $\mathbf b=\mathbf f(\mathbf a)$. Find the local inverse $\mathbf g$ with $\mathbf g(\mathbf b)=\mathbf a$, calculate $D\mathbf f(\mathbf a)$ and $D\mathbf g(\mathbf b)$, and verify formula (52) of the inverse function theorem.
\item Describe the images of lines parallel to the coordinate axes.
\end{exercisesubparts}
\end{exercise}
\begin{solution}
\begin{exercisesubparts}
\item Since $e^x>0$ and $y$ is an arbitrary angle,
$\mathbf f(\R^2)=\R^2\setminus\{\mathbf0\}$.
\item The Jacobian matrix is
\[
 D\mathbf f(x,y)=e^x
 \begin{pmatrix}\cos y&-\sin y\\\sin y&\cos y\end{pmatrix},
 \qquad \det D\mathbf f(x,y)=e^{2x}>0.
\]
The inverse function theorem gives local injectivity. Globally $\mathbf f(x,y+2\pi)=\mathbf f(x,y)$, so the map is not injective.
\item We have $\mathbf b=(1/2,\sqrt3/2)$. On a neighborhood with $u>0$ and $v>0$, choose the angle near $\pi/3$ and set
\[
 \mathbf g(u,v)=\left(\frac12\log(u^2+v^2),\arctan\frac vu\right).
\]
A direct computation yields
\[
 D\mathbf f(\mathbf a)=
 \begin{pmatrix}1/2&-\sqrt3/2\\\sqrt3/2&1/2\end{pmatrix},\qquad
 D\mathbf g(\mathbf b)=
 \begin{pmatrix}1/2&\sqrt3/2\\-\sqrt3/2&1/2\end{pmatrix}.
\]
Their product is $\mathbf I_2$, which verifies $D\mathbf g(\mathbf b)=[D\mathbf f(\mathbf a)]^{-1}$.
\item The lines $x=c$ map onto circles centered at the origin of radius $e^c$. The lines $y=c$ map onto open rays $\{r(\cos c,\sin c):r>0\}$.
\end{exercisesubparts}
\end{solution}

\begin{exercise}{9.18}
Answer questions analogous to Exercise 9.17 for
\[
 \mathbf f(x,y)=(u,v)=(x^2-y^2,2xy).
\]
\end{exercise}
\begin{solution}
This map is the complex squaring map $z\mapsto z^2$, with $z=x+iy$ and $u+iv=z^2$.
\begin{exercisesubparts}
\item Every complex number has a square root; hence the range is all of $\R^2$.
\item
\[
 D\mathbf f(x,y)=\begin{pmatrix}2x&-2y\\2y&2x\end{pmatrix},
 \qquad \det D\mathbf f(x,y)=4(x^2+y^2).
\]
Thus $\mathbf f$ has a local inverse at each nonzero point but not at the origin by the inverse function theorem. Also $\mathbf f(-x,-y)=\mathbf f(x,y)$, so it is not injective. At the origin it cannot be locally injective, since every neighborhood contains $\mathbf z$ and $-\mathbf z$.
\item Taking, as in Exercise 9.17, $\mathbf a=(0,c)$ with $c=\pi/3$, we obtain $\mathbf b=(-c^2,0)$. Write $r=\sqrt{u^2+v^2}$. On a neighborhood of $\mathbf b$ with $r-u>0$, the branch of the square root with positive imaginary part is
\[
 \mathbf g(u,v)=\left(
   \frac{v}{\sqrt{2(r-u)}},\sqrt{\frac{r-u}{2}}
 \right).
\]
It satisfies $\mathbf g(\mathbf b)=\mathbf a$. Moreover,
\[
 D\mathbf f(\mathbf a)=\begin{pmatrix}0&-2c\\2c&0\end{pmatrix},
 \quad D\mathbf g(\mathbf b)=\begin{pmatrix}0&1/(2c)\\-1/(2c)&0\end{pmatrix},
\]
so $D\mathbf g(\mathbf b)=[D\mathbf f(\mathbf a)]^{-1}$.
\item The image of a horizontal line $y=c$ is the parabola
\[
 (u,v)=(x^2-c^2,2cx)\qquad(x\in\R).
\]
For $c\ne0$ it has equation $u=v^2/(4c^2)-c^2$; for $c=0$ it is the nonnegative $u$-axis. Similarly, the image of a vertical line $x=c$ is $u=c^2-v^2/(4c^2)$ for $c\ne0$, and the nonpositive $u$-axis for $c=0$.
\end{exercisesubparts}
\end{solution}

\begin{exercise}{9.19}
Show that the system
\[
 \begin{aligned}
 3x+y-z+u^2&=0,\\
 x-y+2z+u&=0,\\
 2x+2y-3z+2u&=0
 \end{aligned}
\]
can be solved for $x,y,u$ in terms of $z$; for $x,z,u$ in terms of $y$; for $y,z,u$ in terms of $x$; but not for $x,y,z$ in terms of $u$.
\end{exercise}
\begin{solution}
The last two equations and the first together reduce, by elimination, to
\[
 4x+z+4u=0,\qquad 4y-7z=0,\qquad u(u-3)=0.
\]
Consequently all solutions lie on two affine lines:
\[
 \begin{array}{c|ccc}
 &x&y&u\\\hline
 \text{first branch}&-z/4&7z/4&0\\
 \text{second branch}&-z/4-3&7z/4&3
 \end{array}.
\]
On either branch, $z$ parametrizes the solutions. Since $x=-z/4-u$ and $y=7z/4$, we may also use $x$ or $y$ as a parameter and solve uniquely for the other variables in a neighborhood of any fixed solution.

This agrees with the implicit function theorem: if $\mathbf F$ is the left-hand side, the relevant $3\times3$ Jacobian minors are
\[
 \det D_{(x,y,u)}\mathbf F=4(2u-3),\quad
 \det D_{(x,z,u)}\mathbf F=-7(2u-3),\quad
 \det D_{(y,z,u)}\mathbf F=3-2u,
\]
which are nonzero at $u=0$ and $u=3$. In contrast,
$\det D_{(x,y,z)}\mathbf F=0$ everywhere. More directly, fixing a permitted value $u\in\{0,3\}$ leaves the parameter $z$ arbitrary, so $x,y,z$ cannot be uniquely determined as functions of $u$.
\end{solution}

\begin{exercise}{9.20}
Take $n=m=1$ in the implicit function theorem and interpret both the theorem and its proof graphically.
\end{exercise}
\begin{solution}
Let $F\in C^1$ near $(a,b)\in\R^2$, with $F(a,b)=0$ and $\partial_yF(a,b)\ne0$. Then there exist neighborhoods of $a$ and $b$ and a unique $C^1$ function $g$ near $a$, with $g(a)=b$, such that the nearby zero set of $F$ is the graph
\[
 y=g(x),\qquad F(x,g(x))=0.
\]
Differentiating along that graph gives
\[
 g'(x)=-\frac{\partial_xF(x,g(x))}{\partial_yF(x,g(x))}.
\]
Geometrically, $\partial_yF\ne0$ means that the level curve $F=0$ has a nonvertical tangent and can be represented locally as a graph over the $x$-axis. In the proof, the map
\[
 H(x,y)=(x,F(x,y))
\]
has $\det DH(a,b)=\partial_yF(a,b)\ne0$. The inverse function theorem supplies local coordinates $(x,v)$ with $v=F(x,y)$; the level curve becomes the straight line $v=0$. Applying $H^{-1}$ to that line produces the graph $y=g(x)$.
\end{solution}

\begin{exercise}{9.21}
Define $f\colon\R^2\to\R$ by
\[
 f(x,y)=2x^3-3x^2+2y^3+3y^2.
\]
\begin{exercisesubparts}
\item Find the four points where $\nabla f=\mathbf0$, and show that $f$ has exactly one local maximum and one local minimum in $\R^2$.
\item Let $S=\{(x,y)\in\R^2:f(x,y)=0\}$. Find the points of $S$ having no neighborhood in which the equation $f(x,y)=0$ can be solved for $y$ in terms of $x$ (or for $x$ in terms of $y$). Describe $S$ as precisely as possible.
\end{exercisesubparts}
\end{exercise}
\begin{solution}
\begin{exercisesubparts}
\item Since
\[
 \nabla f=(6x(x-1),6y(y+1)),\qquad
 \Hess f=\begin{pmatrix}12x-6&0\\0&12y+6\end{pmatrix},
\]
the critical points are $(0,0)$, $(1,0)$, $(0,-1)$, and $(1,-1)$. The Hessian is positive definite at $(1,0)$ and negative definite at $(0,-1)$, giving a strict local minimum with value $-1$ and a strict local maximum with value $1$, respectively. The Hessian is indefinite at $(0,0)$ and $(1,-1)$, so these are saddle points. Every local extremum of a differentiable function is critical (Exercise 9.8), hence there are no others.
\item The polynomial factors as
\[
 f(x,y)=(x+y)(2x^2-2xy+2y^2-3x+3y).
\]
Writing $X=x-1/2$ and $Y=y+1/2$, the second factor vanishes precisely on the ellipse
\[
 X^2-XY+Y^2=\frac34.
\]
Thus $S$ is the union of this ellipse, centered at $(1/2,-1/2)$, and the line $y=-x$. They meet transversely at the two saddle points $(0,0)$ and $(1,-1)$; neither coordinate can be expressed as a single-valued function of the other on the whole nearby zero set at those crossings.

At every other point $S$ is a smooth curve, and the implicit function theorem applies whenever the appropriate partial derivative is nonzero. On the ellipse, the points where $\partial_y f=0$ (vertical tangents) are
\[
 \left(-\frac12,-1\right),\qquad \left(\frac32,0\right).
\]
At these two points $y$ cannot locally be solved as a single-valued function of $x$; this is seen directly from the strict local extrema of the ellipse's $x$-coordinate. Likewise, the horizontal tangents occur at
\[
 \left(0,-\frac32\right),\qquad \left(1,\frac12\right),
\]
where $x$ cannot be solved as a function of $y$. Therefore the failures for $y=y(x)$ are the two crossings together with $(-1/2,-1)$ and $(3/2,0)$; the failures for $x=x(y)$ are the two crossings together with $(0,-3/2)$ and $(1,1/2)$.
\end{exercisesubparts}
\end{solution}

\begin{exercise}{9.22}
Give a discussion analogous to Exercise 9.21 for
\[
 f(x,y)=2x^3+6xy^2-3x^2+3y^2.
\]
\end{exercise}
\begin{solution}
The first derivatives and Hessian are
\[
 \nabla f=(6(x^2+y^2-x),6y(2x+1)),\qquad
 \Hess f=\begin{pmatrix}12x-6&12y\\12y&12x+6\end{pmatrix}.
\]
Solving $\nabla f=\mathbf0$ gives only $(0,0)$ and $(1,0)$: when $y=0$, the first equation yields $x=0$ or $1$, while $x=-1/2$ is incompatible with $x^2+y^2=x$. The Hessian at $(0,0)$ is indefinite, so $(0,0)$ is a saddle point; at $(1,0)$ it is positive definite, so $(1,0)$ is a strict local minimum with $f(1,0)=-1$. There is no local maximum, since there are no further critical points.

The zero set $S$ is determined by
\[
 3(2x+1)y^2=x^2(3-2x).
\]
There are no zeros with $x\le-1/2$ or $x>3/2$, and for $-1/2<x\le3/2$ the two branches are
\[
 y=\pm |x|\sqrt{\frac{3-2x}{3(2x+1)}}.
\]
They cross at the node $(0,0)$ (with tangent slopes $\pm1$), join at $(3/2,0)$, and extend to $|y|\to\infty$ as $x\downarrow-1/2$. At the node the zero set cannot be represented by a single-valued graph in either coordinate.

For $y$ as a function of $x$, the only further obstruction is the vertical tangent at $(3/2,0)$, where $\partial_yf=0$ but $\partial_xf\ne0$; the curve folds in its $x$-projection.
For $x$ as a function of $y$, further obstructions occur at the two horizontal tangents where $\partial_xf=0$ and $\partial_yf\ne0$. Substitution in the zero-set equation gives
\[
 \left(\frac{\sqrt3}{2},\ \frac12\sqrt{2\sqrt3-3}\right),\qquad
 \left(\frac{\sqrt3}{2},-\frac12\sqrt{2\sqrt3-3}\right).
\]
These are genuine folds: $\partial_{xx}f=6(\sqrt3-1)\ne0$ there. At every remaining point of $S$, the respective implicit function is locally well defined.
\end{solution}

\begin{exercise}{9.23}
Define $f\colon\R^3\to\R$ by
\[
 f(x,y_1,y_2)=x^2y_1+e^x+y_2.
\]
Show that $f(0,1,-1)=0$, that $\partial_1f(0,1,-1)\ne0$, and hence that there is a differentiable function $g$ near $(1,-1)$ in $\R^2$ such that $g(1,-1)=0$ and
\[
 f(g(y_1,y_2),y_1,y_2)=0.
\]
Find $\partial_1g(1,-1)$ and $\partial_2g(1,-1)$.
\end{exercise}
\begin{solution}
We have $f(0,1,-1)=0$ and
\[
 \partial_xf=2xy_1+e^x,\qquad
 \partial_{y_1}f=x^2,\qquad \partial_{y_2}f=1.
\]
Thus $\partial_xf(0,1,-1)=1\ne0$, so the implicit function theorem yields the required $g$. Differentiate the identity $f(g(y_1,y_2),y_1,y_2)=0$:
\[
 \partial_jg=-\frac{\partial_{y_j}f}{\partial_xf}
 \quad(j=1,2),
\]
with derivatives of $f$ evaluated at $(g(y_1,y_2),y_1,y_2)$. At $(1,-1)$, this gives
\[
 \boxed{\partial_1g(1,-1)=0,\qquad\partial_2g(1,-1)=-1.}
\]
\end{solution}

\begin{exercise}{9.24}
For $(x,y)\ne(0,0)$, define $\mathbf f=(f_1,f_2)$ by
\[
 f_1(x,y)=\frac{x^2-y^2}{x^2+y^2},\qquad
 f_2(x,y)=\frac{xy}{x^2+y^2}.
\]
Compute the rank of $D\mathbf f(x,y)$, and find the range of $\mathbf f$.
\end{exercise}
\begin{solution}
Write $(x,y)=r(\cos\theta,\sin\theta)$ with $r>0$. Then
\[
 \mathbf f(x,y)=\left(\cos2\theta,\frac12\sin2\theta\right).
\]
Hence the range is precisely the ellipse
\[
 \boxed{\{(u,v)\in\R^2:u^2+4v^2=1\}.}
\]
The radial derivative of $\mathbf f$ is zero, since $\mathbf f$ is constant on each ray. The angular derivative is $(-2\sin2\theta,\cos2\theta)$, which never vanishes. As the radial and angular directions span $\R^2$ at every nonzero point, the derivative has exactly one-dimensional range:
\[
 \boxed{\rk D\mathbf f(x,y)=1\qquad((x,y)\ne(0,0)).}
\]
\end{solution}

\begin{exercise}{9.25}
Suppose $\mathbf A\in\mathcal L(\R^n,\R^m)$ has rank $r$.
\begin{exercisesubparts}
\item Define $\mathbf S$ as in the proof of Theorem 9.32. Prove that $\mathbf S\mathbf A$ is a projection of $\R^n$ whose null space is $\ker\mathbf A$ and whose range is $\im\mathbf S$. 
\textit{Hint.} By formula (68), $\mathbf S\mathbf A\mathbf S\mathbf A=\mathbf S\mathbf A$.
\item Use (a) to prove $\dim\ker\mathbf A+\dim\im\mathbf A=n$.
\end{exercisesubparts}
\end{exercise}
\begin{solution}
Choose a basis $\mathbf y_1,\ldots,\mathbf y_r$ for $\im\mathbf A$ and vectors $\mathbf z_i\in\R^n$ with $\mathbf A\mathbf z_i=\mathbf y_i$. As in the proof of Theorem 9.32, define $\mathbf S\colon\im\mathbf A\to\R^n$ by
\[
 \mathbf S\left(\sum_{i=1}^r c_i\mathbf y_i\right)=\sum_{i=1}^r c_i\mathbf z_i.
\]
Then $\mathbf A\mathbf S\mathbf y=\mathbf y$ for every $\mathbf y\in\im\mathbf A$.
\begin{exercisesubparts}
\item Put $\mathbf P=\mathbf S\mathbf A$. Since $\mathbf A\mathbf S$ is the identity on $\im\mathbf A$,
\[
 \mathbf P^2=\mathbf S(\mathbf A\mathbf S)\mathbf A=\mathbf S\mathbf A=\mathbf P,
\]
so $\mathbf P$ is a projection. If $\mathbf A\mathbf x=\mathbf0$, then $\mathbf P\mathbf x=\mathbf0$. Conversely, if $\mathbf P\mathbf x=\mathbf0$, then
$\mathbf A\mathbf x=\mathbf A\mathbf S\mathbf A\mathbf x=\mathbf A\mathbf P\mathbf x=\mathbf0$.
Thus $\ker\mathbf P=\ker\mathbf A$. Clearly $\im\mathbf P\subseteq\im\mathbf S$; for $\mathbf y\in\im\mathbf A$, the equality $\mathbf P(\mathbf S\mathbf y)=\mathbf S\mathbf y$ proves the reverse inclusion.
\item Every projection decomposes its space as the direct sum of its null space and range (Section 9.31):
\[
 \R^n=\ker\mathbf P\oplus\im\mathbf P.
\]
The map $\mathbf S$ is injective because $\mathbf A\mathbf S$ is the identity on $\im\mathbf A$, so $\dim\im\mathbf S=\dim\im\mathbf A=r$. Therefore
\[
 \boxed{\dim\ker\mathbf A+\dim\im\mathbf A=n.}
\]
\end{exercisesubparts}
\end{solution}

\begin{exercise}{9.26}
Show that the existence (and even continuity) of $\partial_1\partial_2f$ does not imply the existence of $\partial_1f$. For instance, take $f(x,y)=g(x)$, where $g$ is nowhere differentiable.
\end{exercise}
\begin{solution}
Since $f$ is independent of $y$, we have $\partial_2f(x,y)=0$ for every $(x,y)\in\R^2$. Therefore $\partial_1\partial_2f=0$ everywhere, and this mixed partial is continuous. On the other hand, the difference quotient for $\partial_1f(x,y)$ is exactly the difference quotient for $g'(x)$; it has no limit at any point by the choice of $g$.
\end{solution}

\begin{exercise}{9.27}
Define $f(0,0)=0$ and
\[
 f(x,y)=\frac{xy(x^2-y^2)}{x^2+y^2}\qquad((x,y)\ne(0,0)).
\]
Prove that
\begin{exercisesubparts}
\item $f,\partial_1f,\partial_2f$ are continuous on $\R^2$;
\item $\partial_1\partial_2f$ and $\partial_2\partial_1f$ exist everywhere and are continuous except at the origin;
\item $(\partial_1\partial_2f)(0,0)=1$ and $(\partial_2\partial_1f)(0,0)=-1$.
\end{exercisesubparts}
\end{exercise}
\begin{solution}
\begin{exercisesubparts}
\item The bound
\[
 |f(x,y)|\le|xy|\le\frac12(x^2+y^2)
\]
proves continuity of $f$ at the origin. Its two first partial derivatives there vanish, because $f(h,0)=f(0,h)=0$. Away from the origin,
\[
 \partial_1f=\frac{y(x^4+4x^2y^2-y^4)}{(x^2+y^2)^2},\qquad
 \partial_2f=\frac{x(x^4-4x^2y^2-y^4)}{(x^2+y^2)^2}.
\]
Both are homogeneous of degree one and bounded on the unit circle by continuity; hence both tend to zero as $(x,y)\to(0,0)$. Thus $f\in C^1(\R^2)$.
\item For $(x,y)\ne(0,0)$, rational differentiation gives
\[
 \partial_1\partial_2f
 =\partial_2\partial_1f
 =\frac{(x^2-y^2)(x^4+10x^2y^2+y^4)}{(x^2+y^2)^3}.
\]
These functions are continuous away from the origin. Their values on the $x$-axis are $1$, whereas on the $y$-axis they are $-1$. Thus neither is continuous at $(0,0)$, provided the values there exist as shown next.
\item Using the formulas in (a), for every nonzero $x,y$ one has
\[
 (\partial_2f)(x,0)=x,\qquad
 (\partial_1f)(0,y)=-y.
\]
Both first partials vanish at the origin. Hence by their defining difference quotients,
\[
 (\partial_1\partial_2f)(0,0)=\lim_{h\to0}\frac{h}{h}=1,
\quad
 (\partial_2\partial_1f)(0,0)=\lim_{h\to0}\frac{-h}{h}=-1.
\]
This shows why continuity is essential in the mixed-partials theorem (Theorem 9.41).
\end{exercisesubparts}
\end{solution}

\begin{exercise}{9.28}
For $t\ge0$, put
\[
 \varphi(x,t)=
 \begin{cases}
 x,&0\le x\le\sqrt t,\\
 -x+2\sqrt t,&\sqrt t\le x\le2\sqrt t,\\
 0,&\text{otherwise},
 \end{cases}
\]
and define $\varphi(x,t)=-\varphi(x,|t|)$ for $t<0$. Show that $\varphi$ is continuous on $\R^2$ and that $\partial_2\varphi(x,0)=0$ for every $x$. Define
\[
 F(t)=\int_{-1}^1\varphi(x,t)\,\dd x.
\]
Show that $F(t)=t$ if $|t|<1/4$, so
\[
 F'(0)\ne\int_{-1}^1(\partial_2\varphi)(x,0)\,\dd x.
\]
\end{exercise}
\begin{solution}
For each fixed $t\ne0$, the two linear pieces meet at height $\sqrt{|t|}$ and vanish at the endpoints of their support; hence $\varphi$ is continuous where $t\ne0$. Moreover,
\[
 |\varphi(x,t)|\le\sqrt{|t|}
\]
for all $x,t$. Since $\varphi(x,0)=0$, this proves continuity also at every point of the line $t=0$.

For fixed $x\ne0$, the support $[0,2\sqrt{|t|}]$ eventually misses $x$ as $t\to0$, so $\varphi(x,t)=0$ for all sufficiently small $|t|$. At $x=0$, $\varphi(0,t)=0$ identically. Therefore $\partial_2\varphi(x,0)=0$ for every $x$.

If $0<t<1/4$, the support is contained in $[-1,1]$, and the graph of $\varphi(\cdot,t)$ is a triangle of base $2\sqrt t$ and height $\sqrt t$. Its integral is $t$. By the definition for negative $t$, the integral is $t$ also for $-1/4<t<0$. Hence $F(t)=t$ near $0$ and
\[
 \boxed{F'(0)=1\ne0=\int_{-1}^1\partial_2\varphi(x,0)\,\dd x.}
\]
Continuity of $\varphi$ alone does not justify differentiating under the integral sign.
\end{solution}

\begin{exercise}{9.29}
Let $E\subseteq\R^n$ be open and $f\in C^k(E)$ for a positive integer $k$. Show, by repeated application of Theorem 9.41, that the $k$th-order partial derivative
\[
 \partial_{i_1}\partial_{i_2}\cdots\partial_{i_k}f
\]
is unchanged by permuting the indices $i_1,\ldots,i_k$. For example, if $n\ge3$ and $f\in C^4(E)$, then $\partial_1\partial_2\partial_1\partial_3 f=\partial_3\partial_1\partial_1\partial_2 f$.
\end{exercise}
\begin{solution}
Every permutation of $k$ indices is a product of transpositions of adjacent indices. It is therefore enough to interchange two consecutive partial derivatives in an iterated derivative. Suppose the adjacent indices are $i$ and $j$, and let $h$ be the result of applying to $f$ all partial derivatives that occur to their right. Since $f\in C^k(E)$, the function $h$ has continuous second-order partial derivatives. Theorem 9.41 gives
\[
 \partial_i\partial_jh=\partial_j\partial_ih.
\]
Applying any derivatives that occur to their left preserves the equality. Thus each adjacent transposition leaves the $k$th derivative unchanged, and so does every permutation. The displayed fourth-order identity is one instance.
\end{solution}

\begin{exercise}{9.30}
Let $f\in C^m(E)$, where $E\subseteq\R^n$ is open. Fix $\mathbf a\in E$ and choose $\mathbf x\in\R^n$ sufficiently small that $\mathbf p(t)=\mathbf a+t\mathbf x\in E$ for $0\le t\le1$. Define $h(t)=f(\mathbf p(t))$ wherever $\mathbf p(t)\in E$.
\begin{exercisesubparts}
\item For $1\le k\le m$, show that
\[
 h^{(k)}(t)=\sum_{i_1,\ldots,i_k=1}^n
 (\partial_{i_1}\cdots\partial_{i_k}f)(\mathbf p(t))
 x_{i_1}\cdots x_{i_k}.
\]
The sum is over all ordered $k$-tuples of indices.
\item Taylor's theorem (Theorem 5.15) gives $h(1)=\sum_{k=0}^{m-1}h^{(k)}(0)/k!+h^{(m)}(\theta)/m!$ for some $\theta\in(0,1)$. Use it to show that
\[
 f(\mathbf a+\mathbf x)
 =\sum_{k=0}^{m-1}\frac1{k!}
 \sum_{i_1,\ldots,i_k=1}^n
 (\partial_{i_1}\cdots\partial_{i_k}f)(\mathbf a)
 x_{i_1}\cdots x_{i_k}+r(\mathbf x),
\]
where the $k=0$ term is $f(\mathbf a)$ and
$\lim_{\mathbf x\to\mathbf0}r(\mathbf x)/\|\mathbf x\|^{m-1}=0$.
\item By Exercise 9.29, different orders of a fixed collection of indices give the same derivative. For example, $\partial_1\partial_1\partial_3f$ occurs in three ordered forms. Group the repetitions to show that the Taylor polynomial in (b) equals
\[
 \sum_{\substack{s_1,\ldots,s_n\in\Nzero\\s_1+\cdots+s_n\le m-1}}
 \frac{(\partial_1^{s_1}\cdots\partial_n^{s_n}f)(\mathbf a)}{s_1!\cdots s_n!}
 x_1^{s_1}\cdots x_n^{s_n}.
\]
\end{exercisesubparts}
\end{exercise}
\begin{solution}
\begin{exercisesubparts}
\item The chain rule gives $h'(t)=\sum_i(\partial_if)(\mathbf p(t))x_i$. Differentiating once more applies the chain rule to each partial derivative, producing one additional index and factor $x_i$. Induction gives the displayed formula for every $k\le m$.
\item The one-variable Taylor theorem gives, for some $\theta\in(0,1)$,
\[
 h(1)=\sum_{k=0}^{m-1}\frac{h^{(k)}(0)}{k!}
       +\frac{h^{(m)}(\theta)}{m!}.
\]
Substitute (a) to obtain the requested polynomial. Since the $m$th-order partial derivatives of $f$ are continuous near $\mathbf a$, they are bounded on a sufficiently small closed ball about $\mathbf a$, say by $M$. Thus
\[
 |r(\mathbf x)|\le\frac M{m!}\sum_{i_1,\ldots,i_m=1}^n
 |x_{i_1}\cdots x_{i_m}|
 =\frac M{m!}\left(\sum_{i=1}^n|x_i|\right)^m
 \le\frac{Mn^{m/2}}{m!}\|\mathbf x\|^m.
\]
Consequently $r(\mathbf x)/\|\mathbf x\|^{m-1}\to0$.
\item By Exercise 9.29, all ordered derivative strings containing $s_i$ occurrences of index $i$ have the same value. If $s_1+\cdots+s_n=k$, there are exactly
$k!/(s_1!\cdots s_n!)$ such strings. Their combined coefficient after multiplication by $1/k!$ is therefore $1/(s_1!\cdots s_n!)$. Summing over all nonnegative multi-indices of total degree at most $m-1$ yields the stated polynomial.
\end{exercisesubparts}
\end{solution}

\begin{exercise}{9.31}
Suppose $f\in C^3$ near $\mathbf a\in\R^2$, $\nabla f(\mathbf a)=\mathbf0$, and not all second-order partial derivatives vanish at $\mathbf a$. Show how the degree-two Taylor polynomial at $\mathbf a$ can be used to decide whether $f$ has a local maximum, a local minimum, or neither. Extend the discussion to $\R^n$.
\end{exercise}
\begin{solution}
Put $\mathbf h=\mathbf x-\mathbf a$. By Exercise 9.30 and $\nabla f(\mathbf a)=\mathbf0$,
\[
 f(\mathbf a+\mathbf h)-f(\mathbf a)=Q(\mathbf h)+o(\|\mathbf h\|^2),
 \qquad Q(\mathbf h)=\frac12\mathbf h^{\mathrm T}(\Hess f(\mathbf a))\mathbf h.
\]
If $Q$ is positive definite, compactness of the unit sphere gives $Q(\mathbf h)\ge c\|\mathbf h\|^2$ for some $c>0$, so $\mathbf a$ is a strict local minimum. If $Q$ is negative definite, it is a strict local maximum. If $Q$ takes both positive and negative values, then along two suitable lines the quadratic term determines opposite signs for sufficiently small nonzero $\mathbf h$, so $\mathbf a$ is neither a local maximum nor a local minimum.

In two variables, write
\[
 \Hess f(\mathbf a)=\begin{pmatrix}A&B\\B&C\end{pmatrix},
 \qquad\Delta=AC-B^2.
\]
If $\Delta>0$, the form is definite: $A>0$ gives a strict minimum and $A<0$ gives a strict maximum. If $\Delta<0$, the form is indefinite and $\mathbf a$ is a saddle point. If $\Delta=0$, the nonzero quadratic form is semidefinite, and the Taylor polynomial of degree two does \emph{not} always decide the question.

Indeed, $f_+(x,y)=x^2+y^4$ and $f_-(x,y)=x^2-y^4$ belong to $C^3$, have zero gradient at the origin, and have the same nonzero quadratic Taylor polynomial $Q(x,y)=x^2$. Nevertheless, $f_+$ has a strict local minimum and $f_-$ has neither a local minimum nor a local maximum at the origin. Thus the exercise as worded needs this qualification.

The same reasoning holds in $\R^n$: classify the quadratic form associated with the symmetric Hessian. Positive definite implies a strict local minimum, negative definite a strict local maximum, and indefinite implies neither. A nonzero semidefinite Hessian leaves the behavior along its null directions undecided; higher-order terms must then be examined.
\end{solution}


\begin{exercise}{11.1}
If $f\ge0$ and $\int_E f\,\dd\mu=0$, prove that $f(x)=0$ almost everywhere on $E$.

\textit{Hint.} Let $E_n$ be the subset of $E$ on which $f(x)>1/n$. Write $A=\bigcup_{n=1}^\infty E_n$. Then $\mu(A)=0$ if and only if $\mu(E_n)=0$ for every $n$.
\end{exercise}
\begin{solution}
Put $E_n=\{x\in E:f(x)>1/n\}$. Since $f\ge0$,
\[
 0=\int_E f\,\dd\mu\ge\int_{E_n}f\,\dd\mu
 \ge\frac1n\mu(E_n).
\]
Hence $\mu(E_n)=0$ for every $n$. The set on which $f$ is positive is
\[
 \{x\in E:f(x)>0\}=\bigcup_{n=1}^\infty E_n,
\]
which has measure zero by countable subadditivity. Thus $f=0$ almost everywhere on $E$.
\end{solution}

\begin{exercise}{11.2}
If $\int_A f\,\dd\mu=0$ for every measurable subset $A$ of a measurable set $E$, then $f(x)=0$ almost everywhere on $E$.
\end{exercise}
\begin{solution}
First suppose $f$ is real-valued. The sets
$E_+=\{x\in E:f(x)>0\}$ and $E_-=\{x\in E:f(x)<0\}$ are measurable. By hypothesis,
\[
 \int_{E_+}f\,\dd\mu=0,
 \qquad\int_{E_-}(-f)\,\dd\mu=0.
\]
Exercise 11.1 applied to $f$ on $E_+$ and to $-f$ on $E_-$ shows that both $E_+$ and $E_-$ have measure zero. Hence $f=0$ almost everywhere on $E$.

For complex-valued $f$, the hypothesis implies that the integrals of $\operatorname{Re}f$ and $\operatorname{Im}f$ vanish over every measurable $A\subseteq E$. Applying the real case to these two functions gives the same conclusion.
\end{solution}

\begin{exercise}{11.3}
If $(f_n)$ is a sequence of measurable functions, prove that the set of points $x$ at which $(f_n(x))$ converges is measurable.
\end{exercise}
\begin{solution}
By the Cauchy criterion in $\R$ (or $\C$), the set in question is
\[
 C=\bigcap_{k=1}^\infty\bigcup_{N=1}^\infty
       \bigcap_{m,n\ge N}
       \left\{x:|f_n(x)-f_m(x)|<\frac1k\right\}.
\]
Each set inside the braces is measurable, since $f_n-f_m$ is measurable. Countable unions and intersections of measurable sets are measurable. Thus $C$ is measurable.
\end{solution}

\begin{exercise}{11.4}
If $f\in L^1(\mu)$ on $E$ and $g$ is bounded and measurable on $E$, then $fg\in L^1(\mu)$ on $E$.
\end{exercise}
\begin{solution}
The product $fg$ is measurable. Choose $M<\infty$ such that $|g(x)|\le M$ on $E$. Then
\[
 \int_E|fg|\,\dd\mu\le M\int_E|f|\,\dd\mu<\infty.
\]
Hence $fg$ is Lebesgue integrable on $E$.
\end{solution}

\begin{exercise}{11.5}
Put
\[
 g(x)=\begin{cases}0,&0\le x\le\frac12,\\1,&\frac12<x\le1,\end{cases}
\qquad
 f_{2k}(x)=g(x),\quad f_{2k+1}(x)=g(1-x)
 \qquad(0\le x\le1).
\]
Show that
\[
 \liminf_{n\to\infty}f_n(x)=0\qquad(0\le x\le1),
\]
but
\[
 \int_0^1 f_n(x)\,\dd x=\frac12.
\]
\textit{Compare with (77).}
\end{exercise}
\begin{solution}
At each $x\in[0,1]$, at least one of $g(x)$ and $g(1-x)$ is zero. Since the sequence alternates between these two functions, its lower limit is zero at every point. On the other hand, $g$ and $g(1-\cdot)$ are characteristic functions of intervals of length $1/2$, up to endpoint values. Hence, for every $n$,
\[
 \int_0^1 f_n(x)\,\dd x=\frac12.
\]
In particular, the lower limit of the functions does not determine the limit of their integrals: $\liminf_n\int_0^1 f_n\,\dd x=1/2$, whereas $\int_0^1\liminf_n f_n\,\dd x=0$. This is consistent with Fatou's inequality (77).
\end{solution}

\begin{exercise}{11.6}
Let
\[
 f_n(x)=\begin{cases}\frac1n,&|x|\le n,\\0,&|x|>n.\end{cases}
\]
Then $f_n(x)\to0$ uniformly on $\R$, but
\[
 \int_{-\infty}^{\infty}f_n(x)\,\dd x=2
 \qquad(n=1,2,3,\ldots).
\]
Thus uniform convergence does not imply dominated convergence in the sense of Theorem 11.32. However, on sets of finite measure, uniformly convergent sequences of bounded functions do satisfy Theorem 11.32.
\end{exercise}
\begin{solution}
For every $n$, $\sup_{x\in\R}|f_n(x)|=1/n$, which tends to zero. Nevertheless,
\[
 \int_{\R}f_n\,\dd x=\frac1n\,m([-n,n])=2.
\]
Consequently there can be no integrable function dominating all $f_n$, since such a function would yield $\int f_n\to0$ by Theorem 11.32.

Now let $E$ have finite measure and suppose that measurable, bounded functions $f_n$ converge uniformly to $f$ on $E$. The limit $f$ is measurable, and
\[
 \int_E|f_n-f|\,\dd\mu
 \le\mu(E)\sup_{x\in E}|f_n(x)-f(x)|\longrightarrow0.
\]
In particular, their integrals converge. Moreover, uniform convergence and boundedness of the finitely many initial functions give a common bound $|f_n|\le M$ on $E$ for all $n$. Since $M\mathds{1}_E$ is integrable, the hypotheses of the dominated convergence theorem are satisfied.
\end{solution}

\begin{exercise}{11.7}
Find a necessary and sufficient condition that $f\in\mathcal R([a,b],\alpha)$.

\textit{Hint.} Consider Example 11.6(b) and Theorem 11.33.
\end{exercise}
\begin{solution}
Here $\alpha$ is increasing on $[a,b]$, as in Definition 6.2, and $f$ is real and bounded. We must distinguish the two sides of a jump of $\alpha$.

For $c\in(a,b]$, put $j^-(c)=\alpha(c)-\alpha(c-)$, and for $c\in[a,b)$ put $j^+(c)=\alpha(c+)-\alpha(c)$. Removing all these jumps leaves a continuous increasing function $\alpha_0$:
\[
 \alpha_0(x)=\alpha(x)
 -\sum_{c\in(a,b]}j^-(c)\mathds{1}_{[c,b]}(x)
 -\sum_{c\in[a,b)}j^+(c)\mathds{1}_{(c,b]}(x).
\]
The sums are over at most countably many nonzero terms and converge uniformly, since their total is at most $\alpha(b)-\alpha(a)$. Let $\mu_{\alpha_0}$ be the measure constructed from $\alpha_0$ in Example 11.6(b), and let $D_f$ be the set of discontinuities of $f$.

The necessary and sufficient conditions are
\begin{exercisesubparts}
\item $\mu_{\alpha_0}(D_f)=0$;
\item $f(c-)=f(c)$ whenever $j^-(c)>0$, and $f(c+)=f(c)$ whenever $j^+(c)>0$.
\end{exercisesubparts}
To see this, recall that the Riemann--Stieltjes criterion is
\[
 \inf_P\sum_i\bigl(M_i-m_i\bigr)
                   \bigl(\alpha(x_i)-\alpha(x_{i-1})\bigr)=0,
\]
where $M_i,m_i$ are the supremum and infimum of $f$ on $[x_{i-1},x_i]$. For a continuous increasing integrator $\alpha_0$, the proof of Theorem 11.33, with the measure $\mu_{\alpha_0}$ in place of length, shows that this criterion is equivalent to (a). For the single-jump integrator $\mathds{1}_{[c,b]}$, it is equivalent to continuity of $f$ from the left at $c$; for $\mathds{1}_{(c,b]}$, to continuity from the right. This gives (b).

Finally, the oscillation sum is additive in the integrator and nonnegative for every increasing summand. Necessity follows for $\alpha_0$ and for each jump separately. Conversely, choose a partition making the $\alpha_0$ contribution small, refine it to control finitely many jumps, and note that the remaining jumps contribute at most $2\|f\|_\infty$ times their total size, which can be made arbitrarily small. This proves sufficiency.

\textit{Remark.} If $\alpha$ is continuous, (b) disappears and the criterion reduces to $\mu_\alpha(D_f)=0$. Without continuity of $\alpha$, this simpler formulation is not generally necessary: a jump of $f$ on the side where $\alpha$ has no jump may be harmless.
\end{solution}

\begin{exercise}{11.8}
If $f\in\mathcal R([a,b])$ and $F(x)=\int_a^x f(t)\,\dd t$, prove that $F'(x)=f(x)$ almost everywhere on $[a,b]$.
\end{exercise}
\begin{solution}
By Theorem 11.33(b), a Riemann-integrable function is continuous almost everywhere. At a point $x\in(a,b)$ where $f$ is continuous, for sufficiently small $h$ (positive or negative),
\[
 \frac{F(x+h)-F(x)}h-f(x)
 =\frac1h\int_x^{x+h}\bigl(f(t)-f(x)\bigr)\,\dd t.
\]
The absolute value of the right side is at most
$\sup_{t\text{ between }x\text{ and }x+h}|f(t)-f(x)|$, which tends to zero. Thus $F'(x)=f(x)$ at every interior point of continuity of $f$. The exceptional set, together with the endpoints, has measure zero.
\end{solution}

\begin{exercise}{11.9}
Prove that the function $F$ given by (96) is continuous on $[a,b]$.
\end{exercise}
\begin{solution}
In (96), $F(x)=\int_a^x f(t)\,\dd t$ with $f\in L^1([a,b])$. Suppose $x_n\to x$ in $[a,b]$ and put
$h_n(t)=f(t)\mathds{1}_{[a,x_n]}(t)$. Then $h_n(t)\to f(t)\mathds{1}_{[a,x]}(t)$ for every $t\ne x$, hence almost everywhere, and $|h_n(t)|\le|f(t)|$. By the dominated convergence theorem (Theorem 11.32),
\[
 F(x_n)=\int_a^b h_n(t)\,\dd t\longrightarrow
 \int_a^x f(t)\,\dd t=F(x).
\]
Therefore $F$ is continuous on $[a,b]$.
\end{solution}

\begin{exercise}{11.10}
If $\mu(X)<+\infty$ and $f\in L^2(\mu)$ on $X$, prove that $f\in L^1(\mu)$ on $X$. If $\mu(X)=+\infty$, this is false. For instance, if
\[
 f(x)=\frac1{1+|x|},
\]
then $f\in L^2$ on $\R$, but $f\notin L^1$ on $\R$.
\end{exercise}
\begin{solution}
The Schwarz inequality (Theorem 11.35), applied to $|f|$ and the constant function $1$, gives
\[
 \int_X|f|\,\dd\mu
 \le\left(\int_X|f|^2\,\dd\mu\right)^{1/2}
      \mu(X)^{1/2}<\infty.
\]
For the stated example, symmetry gives
\[
 \int_\R|f|^2\,\dd x
 =2\int_0^\infty\frac{\dd x}{(1+x)^2}=2,
 \qquad
 \int_\R|f|\,\dd x
 =2\int_0^\infty\frac{\dd x}{1+x}=\infty.
\]
Thus finiteness of the underlying measure is essential.
\end{solution}

\begin{exercise}{11.11}
If $f,g\in L^1(\mu)$ on $X$, define the distance between $f$ and $g$ by
\[
 d(f,g)=\int_X|f-g|\,\dd\mu.
\]
Prove that $L^1(\mu)$ is a complete metric space.
\end{exercise}
\begin{solution}
As in the discussion following Theorem 11.36, functions equal almost everywhere are identified. Then $d(f,g)=\|f-g\|_1$ satisfies the metric axioms: nonnegativity and symmetry are immediate, Exercise 11.1 gives definiteness, and $|f-h|\le|f-g|+|g-h|$ gives the triangle inequality.

Let $(f_n)$ be Cauchy in $L^1(\mu)$. Choose a subsequence $(f_{n_k})$ such that
\[
 \|f_{n_{k+1}}-f_{n_k}\|_1<2^{-k}
 \qquad(k\ge1).
\]
Set
\[
 H_m=|f_{n_1}|+\sum_{k=1}^{m}|f_{n_{k+1}}-f_{n_k}|.
\]
The functions $H_m$ increase and satisfy
\[
 \int_XH_m\,\dd\mu
 \le\|f_{n_1}\|_1+\sum_{k=1}^\infty2^{-k}<\infty.
\]
By the monotone convergence theorem, $H=\lim_m H_m$ is integrable and finite almost everywhere. Thus the series of differences
$\sum_k(f_{n_{k+1}}-f_{n_k})$ converges absolutely almost everywhere. It defines a measurable limit $f$ of $(f_{n_k})$, with $|f|\le H$ almost everywhere; hence $f\in L^1(\mu)$.

For every $k$, almost everywhere,
\[
 |f-f_{n_k}|\le\sum_{j=k}^\infty|f_{n_{j+1}}-f_{n_j}|.
\]
Integrating (or applying monotone convergence to the tail) yields
\[
 \|f-f_{n_k}\|_1\le\sum_{j=k}^\infty2^{-j}\longrightarrow0.
\]
Since the original sequence is Cauchy, for $n,n_k$ sufficiently large,
\[
 \|f_n-f\|_1\le\|f_n-f_{n_k}\|_1+\|f_{n_k}-f\|_1
\]
is arbitrarily small. Hence $f_n\to f$ in $L^1(\mu)$, proving completeness.
\end{solution}

\begin{exercise}{11.12}
Suppose
\begin{exercisesubparts}
\item $|f(x,y)|\le1$ if $0\le x\le1$ and $0\le y\le1$;
\item for fixed $x$, $f(x,y)$ is a continuous function of $y$;
\item for fixed $y$, $f(x,y)$ is a continuous function of $x$.
\end{exercisesubparts}
Put
\[
 g(x)=\int_0^1f(x,y)\,\dd y\qquad(0\le x\le1).
\]
Is $g$ continuous?
\end{exercise}
\begin{solution}
Yes. Fix $x\in[0,1]$ and let $x_n\to x$. For each $y\in[0,1]$, assumption (c) gives $f(x_n,y)\to f(x,y)$. By (b), each function $y\mapsto f(x_n,y)$ is continuous, hence measurable, and by (a) all are bounded in absolute value by the integrable function $1$ on $[0,1]$. The dominated convergence theorem therefore gives
\[
 g(x_n)=\int_0^1f(x_n,y)\,\dd y
 \longrightarrow\int_0^1f(x,y)\,\dd y=g(x).
\]
Thus $g$ is continuous. Joint continuity of $f$ is not needed.
\end{solution}

\begin{exercise}{11.13}
Consider the functions
\[
 f_n(x)=\sin(nx)\qquad(n=1,2,3,\ldots,\ -\pi\le x\le\pi)
\]
as points of $L^2$. Prove that the set of these points is closed and bounded, but not compact.
\end{exercise}
\begin{solution}
Orthogonality of the sine functions gives
\[
 \|f_n\|_2^2=\int_{-\pi}^\pi\sin^2(nx)\,\dd x=\pi,
\]
and, for $n\ne m$,
\[
 \|f_n-f_m\|_2^2
 =\|f_n\|_2^2+\|f_m\|_2^2
   -2\int_{-\pi}^\pi\sin(nx)\sin(mx)\,\dd x
 =2\pi.
\]
Thus the set is bounded and its distinct points have a fixed positive separation. Any convergent sequence drawn from the set is Cauchy and must therefore be eventually constant; its limit belongs to the set. Hence the set is closed. On the other hand, the sequence $(f_n)$ has no Cauchy subsequence and therefore no convergent subsequence. The set is not compact.
\end{solution}

\begin{exercise}{11.14}
Prove that a complex function $f$ is measurable if and only if $f^{-1}(V)$ is measurable for every open set $V$ in the plane.
\end{exercise}
\begin{solution}
Write $f=u+iv$, with $u,v$ real. If $f$ is measurable, $u$ and $v$ are measurable by definition. Every open set $V\subseteq\C$ is a countable union of open rectangles with rational endpoints. For each such rectangle
\[
 R=(a,b)\times(c,d),
\]
its inverse image under $f$ is
\[
 f^{-1}(R)=u^{-1}((a,b))\cap v^{-1}((c,d)),
\]
which is measurable. Taking countable unions shows that $f^{-1}(V)$ is measurable.

Conversely, for every $t\in\R$, the sets
$V_t=\{z\in\C:\operatorname{Re}z>t\}$ and
$W_t=\{z\in\C:\operatorname{Im}z>t\}$ are open. The hypothesis gives that
\[
 \{x:u(x)>t\}=f^{-1}(V_t),\qquad
 \{x:v(x)>t\}=f^{-1}(W_t)
\]
are measurable. By the level-set criterion, $u$ and $v$ are measurable, and hence so is $f$.
\end{solution}

\begin{exercise}{11.15}
Let $\mathscr R$ be the ring of all elementary subsets of $(0,1]$. If $0<a\le b\le1$, define
\[
 \varphi([a,b])=\varphi([a,b))
 =\varphi((a,b])=\varphi((a,b))=b-a,
\]
but define
\[
 \varphi((0,b))=\varphi((0,b])=1+b
 \qquad(0<b\le1).
\]
Show that this gives an additive set function $\varphi$ on $\mathscr R$, which is not regular and which cannot be extended to a countably additive set function on a $\sigma$-ring.
\end{exercise}
\begin{solution}
Let $m$ denote ordinary length of elementary sets. Every $A\in\mathscr R$ either contains some interval $(0,\varepsilon)$ or is disjoint from one. Define
\[
 \chi(A)=\begin{cases}
 1,&(0,\varepsilon)\subseteq A\text{ for some }\varepsilon>0,\\
 0,&\text{otherwise}.
 \end{cases}
 \qquad\varphi(A)=m(A)+\chi(A).
\]
This agrees with all the prescribed interval values and shows that the definition is independent of a decomposition into elementary intervals. If $A,B\in\mathscr R$ are disjoint, they cannot both contain an interval $(0,\varepsilon)$. Moreover, $A\cup B$ contains such an interval exactly when one of $A,B$ does. Therefore
$\chi(A\cup B)=\chi(A)+\chi(B)$; since $m$ is additive, so is $\varphi$.

The set $A=(0,1]$ satisfies $\varphi(A)=2$. Any closed elementary $F\subseteq A$ is compact and bounded away from $0$, so $\chi(F)=0$ and $\varphi(F)=m(F)\le1$. The inequality $\varphi(A)\le\varphi(F)+\varepsilon$ in Definition 11.5 is impossible for $\varepsilon<1$. Thus $\varphi$ is not regular.

Finally, the disjoint sets
\[
 A_n=\left(\frac1{n+1},\frac1n\right]
 \qquad(n\ge1)
\]
are elementary, and $\bigcup_{n=1}^\infty A_n=(0,1]$. If a countably additive extension existed on a $\sigma$-ring containing $\mathscr R$, it would give
\[
 2=\varphi((0,1])
 =\sum_{n=1}^\infty\varphi(A_n)
 =\sum_{n=1}^\infty\left(\frac1n-\frac1{n+1}\right)=1,
\]
a contradiction.
\end{solution}

\begin{exercise}{11.16}
Suppose $(n_k)$ is an increasing sequence of positive integers and $E$ is the set of all $x\in(-\pi,\pi)$ at which $(\sin(n_kx))$ converges. Prove that $m(E)=0$.

\textit{Hint.} For every $A\subseteq E$,
\[
 \int_A\sin(n_kx)\,\dd x\to0,
\]
and
\[
 2\int_A\sin^2(n_kx)\,\dd x
 =\int_A\bigl(1-\cos(2n_kx)\bigr)\,\dd x
 \to m(A)\qquad(k\to\infty).
\]
\end{exercise}
\begin{solution}
The set $E$ is measurable by Exercise 11.3. Put $h(x)=\lim_{k\to\infty}\sin(n_kx)$ for $x\in E$. For any measurable $A\subseteq E$, the characteristic function $\mathds{1}_A$ belongs to $L^2([-\pi,\pi])$. Since the normalized sine and cosine functions form orthonormal sets, the Bessel inequality shows that
\[
 \int_A\sin(n_kx)\,\dd x\to0,
 \qquad\int_A\cos(2n_kx)\,\dd x\to0,
\]
as $n_k\to\infty$. Hence
\[
 2\int_A\sin^2(n_kx)\,\dd x\to m(A).
\]
On $A$, $\sin(n_kx)\to h(x)$, and both $|\sin(n_kx)|$ and $\sin^2(n_kx)$ are bounded by $1$. By dominated convergence,
\[
 \int_Ah\,\dd x=0,
 \qquad2\int_Ah^2\,\dd x=m(A).
\]
The first identity holds for \emph{every} measurable $A\subseteq E$. Exercise 11.2 therefore gives $h=0$ almost everywhere on $E$. Applying the second identity with $A=E$ now yields $m(E)=2\int_Eh^2\,\dd x=0$.
\end{solution}

\begin{exercise}{11.17}
Suppose $E\subseteq(-\pi,\pi)$, $m(E)>0$, and $\delta>0$. Use the Bessel inequality to prove that there are at most finitely many integers $n$ such that $\sin(nx)\ge\delta$ for all $x\in E$.
\end{exercise}
\begin{solution}
For each positive integer $k$, let
\[
 c_k=\frac1{\sqrt\pi}\int_E\sin(kx)\,\dd x.
\]
Since $(\sin(kx)/\sqrt\pi)_{k\ge1}$ is orthonormal in $L^2([-\pi,\pi])$, Bessel's inequality applied to $\mathds{1}_E$ gives
\[
 \sum_{k=1}^\infty|c_k|^2\le\|\mathds{1}_E\|_2^2=m(E).
\]
If $n\ne0$ satisfies $\sin(nx)\ge\delta$ on $E$, then
\[
 |c_{|n|}|=\frac1{\sqrt\pi}\left|\int_E\sin(nx)\,\dd x\right|
 \ge\frac{\delta\,m(E)}{\sqrt\pi}.
\]
Consequently only finitely many values of $|n|$ can satisfy the inequality, by the square-summability of the coefficients. The integer $n=0$ does not satisfy it because $\delta>0$. Hence there are only finitely many such integers.
\end{solution}

\begin{exercise}{11.18}
Suppose $f\in L^2(\mu)$, $g\in L^2(\mu)$. Prove that
\[
 \left|\int_X f\overline g\,\dd\mu\right|^2
 =\int_X|f|^2\,\dd\mu\int_X|g|^2\,\dd\mu
\]
if and only if there is a constant $c$ such that $g(x)=cf(x)$ almost everywhere.

\textit{Compare Theorem 11.35.}
\end{exercise}
\begin{solution}
Assume first that $\|f\|_2>0$. Put
\[
 c=\frac{\displaystyle\int_X\overline f g\,\dd\mu}{\|f\|_2^2}.
\]
Expanding the nonnegative square gives
\[
 \begin{aligned}
 \|g-cf\|_2^2
 &=\|g\|_2^2
   -\frac{\left|\int_X\overline f g\,\dd\mu\right|^2}{\|f\|_2^2}.
 \end{aligned}
\]
Thus equality in the Schwarz inequality holds precisely when $\|g-cf\|_2=0$, which is equivalent to $g=cf$ almost everywhere by Exercise 11.1. Conversely, if $g=cf$ almost everywhere, direct substitution gives equality.

\textit{Qualification.} The statement as printed omits the case $\|f\|_2=0$. If $f=0$ almost everywhere, equality holds for \emph{every} $g\in L^2(\mu)$, including nonzero functions which cannot be written as $cf$. The exact equality criterion, valid without a nonzero hypothesis, is therefore: either $f=0$ almost everywhere, or $g=cf$ almost everywhere for some constant $c$.
\end{solution}

\end{document}
